\Needspace{18\baselineskip}
\fancyhead[L]{\small GAUSS TO VON ZACH.}
\begin{center}
GAUSS TO VON ZACH.
\end{center}

\begin{center}
Monatliche Correspondenz. Vol. XIV. Pp. 187–192. August 1806.
\end{center}

\begin{flushright}
Brunswick, 8 July 1806.
\end{flushright}

\par\indent --- --- --- Using the Milan observations of the planet \emph{Pallas} (M.\,C., July issue, 1806, \href{https://www.digitale-sammlungen.de/en/view/bsb10538606?page=100}{p.~90}) that you were kind enough to send me, I calculated new elements (IX.) for its orbit that fit \textsc{Oriani}’s observations as closely as possible. From these elements I derived the opposition, which could not have been done as well with the elements in set VIII, which differ too much. The result is:

\begin{center}
1805. Nov. 29.\ \ $11^{\mathrm{u}}\,18^{\mathrm{m}}\,17^{\mathrm{s}}$ mean time at Seeberg.

Longitude:\quad $67^\circ\,20'\,42''9$

Southern geocentric latitude:\quad $54\;\;30\;\;54.9$
\end{center}

\par\indent To prepare an ephemeris for the next apparition of \emph{Pallas}, I have now based my new elements on the three oppositions of 1803, 1804, 1805 and on the observations of 1802. It turned out that these various positions could no longer be represented quite exactly by a pure ellipse, which is undoubtedly due to the perturbations that have so far been neglected. Even so, the discrepancy is small and can hardly exceed one minute. In particular, the observations of 1803 give an orbital inclination $44''$ greater, and those of 1805 an inclination just as much smaller, than the mean used for the elements below. In the observations of 1802 the node appears about $45''$ behind the adopted mean, and in those of 1804 it appears only that far ahead of it. Here are the new elements (IX.) from which the following ephemeris was calculated:
\Needspace{19\baselineskip}
\begin{center}\begingroup
\fontsize{9}{11}\selectfont
\setlength{\tabcolsep}{2pt}\renewcommand{\arraystretch}{1.12}
\begin{tabular}{lr}
\multicolumn{2}{c}{Epoch on the Seeberg meridian} \\
1802. & $143^\circ\,24'\,1''1$ \\
1803. & $221\;32\;56.4$ \\
1804. & $299\;54\;42.5$ \\
1805. & $18\;3\;37.8$ \\
1806. & $96\;12\;33.2$ \\
1807. & $174\;21\;18.5$ \\
\end{tabular}\endgroup\end{center}
\begin{center}\begingroup
\fontsize{9}{11}\selectfont
\setlength{\tabcolsep}{2pt}\renewcommand{\arraystretch}{1.12}
\begin{tabular}{lr}
daily mean tropical motion .\ .\ . & $770''7816$ \\
Longitude of aphelion in 1805 .\ .\ .\ . & $301^\circ\,11'\,31''1$ \\
Longitude of the ascending node in 1805 . & $172\;30\;5.8$ \\
Eccentricity .\ .\ .\ .\ .\ .\ .\ .\ .\ .\ . & $0.2453840$ \\
Logarithm of the mean distance .\ .\ . & $0.4421021$ \\
Inclination of the orbit .\ .\ .\ .\ .\ .\ .\ . & $34^\circ\,37'\,8''1$ \\
\end{tabular}\endgroup\end{center}

\par\indent According to these elements, the next opposition would occur on 4 May 1807 at 21 hours, at longitude $223^\circ\,53'$ and north geocentric latitude $42^\circ\,18'$. Time will tell how accurate this calculation is. I sent my ephemeris to Prof. \textsc{Harding} at once. He will draw a chart from it, and I hope you will soon treat us to it in the M.\,C.
\Needspace{15\baselineskip}
\SourceLeaf{283}%
\par\indent My next work will concern \emph{Juno}. There will be little to improve in my set V of orbital elements. I have already begun this calculation too. For now, here is my comparison of set V with the last three observations Mr.~\textsc{Bessel} made in May:

\begin{center}\begingroup
\fontsize{9}{11}\selectfont
\setlength{\tabcolsep}{2pt}\renewcommand{\arraystretch}{1.12}
\begin{tabular}{l|l|l|l|l}
& \multicolumn{2}{c|}{Calculated} & \multicolumn{2}{c}{Difference in} \\
\shortstack{Mean time\\at Lilienthal} & \shortstack{Right\\ascension} & \shortstack{Northern\\declination} & \shortstack{Right\\ascension} & Declination \\
\hline
1806. May 17.\ $11^{\mathrm{u}}\,47^{\mathrm{m}}\,59^{\mathrm{s}}$ & $165^\circ\;4'\;11''$ & $9^\circ\;30'\;36''$ & $+0'\,51''$ & $+38''$ \\
19.\ $11\;\;1\;\;38$ & $165\;\;16\;\;34$ & $9\;\;29\;\;23$ & $+0\;\;33$ & $-4$ \\
23.\ $11\;\;31\;\;\;2$ & $165\;\;44\;\;57$ & $9\;\;25\;\;13$ & $+1\;\;11$ & $0$ \\
\cline{3-5}
& & Mean & $+52''$ & $+11''$ \\
\end{tabular}\endgroup\end{center}

\par\indent As soon as I have improved the elements of \emph{Juno}, I will at once calculate an ephemeris for it as well and send it to you and to Prof. \textsc{Harding}. Finally, I will do the same work for \emph{Ceres}, as I hope that by then more observations around the time of opposition will have become known than are available now. [The geocentric course of Pallas according to the elements in set IX is given at intervals of three days from 7 December 1806 to 21 September 1807 for 12 hours, midnight at Seeberg, with right ascension and declination in minutes, distance from the Earth to three decimal places and brightness to four.]
